%=====================================================================
\chapter{The Role of Algorithms, and What Efficient Means}\label{ch:role}
%=====================================================================
\locator{1}{Part I --- Analysing an Algorithm}

\begin{outcomes}
By the end of this chapter you will be able to:
\begin{tight}
  \item \textbf{State} what an algorithm is, and \textbf{distinguish} it from a
        program that implements it.
  \item \textbf{Describe} the random-access machine model this course measures
        against, and \textbf{name} three things it deliberately ignores.
  \item \textbf{Explain} why the interesting question about an algorithm is
        asymptotic, and \textbf{compute} how long a given algorithm takes at a
        given input size.
  \item \textbf{Predict} the ratio by which a running time grows when the input
        doubles, from the algorithm's order of growth, and \textbf{check} the
        prediction by measurement.
  \item \textbf{Explain} why a faster computer running a worse algorithm loses,
        and \textbf{find} the input size at which it starts losing.
\end{tight}
\end{outcomes}

\begin{prereq}
Nothing in this book. You need to be able to read a loop and to have written
programs that sort or search --- \emph{Data Structures} is the published
prerequisite and this chapter assumes only its first fortnight. Appendix~A sets
up the compiler, and you will want it from Section~1.4 onward.
\end{prereq}

\begin{keyterms}
algorithm $\bullet$ problem $\bullet$ instance $\bullet$ correctness $\bullet$
random-access machine $\bullet$ RAM model $\bullet$ primitive operation
$\bullet$ running time $\bullet$ input size $\bullet$ order of growth $\bullet$
asymptotic analysis $\bullet$ constant factor $\bullet$ space complexity
$\bullet$ time--space trade-off
\end{keyterms}

\section{Two Programs That Agree on Every Answer}

Here are two ways to decide whether a sorted array of a million names contains
one particular name.

The first walks the array from the start, comparing as it goes, and stops when
it finds the name or runs off the end. The second looks at the middle element,
discards the half that cannot contain the name, and repeats.

Both are correct. Both return the same answer for every input. Run them on the
same machine with the same compiler and the first takes, on a million names,
about half a million comparisons; the second takes twenty. Not twenty thousand
--- twenty.

That gap is the subject of this course, and the first thing to notice about it
is that it has nothing to do with programming. It is not a faster loop, a
better compiler or a tighter data structure. It is a different \emph{idea},
and the difference between the two ideas survives every change of language,
machine and decade.

\begin{definitionbox}[Problem, Instance, Algorithm]
A \term{problem} is a specification: a description of the input it accepts and
the output it must produce. ``Sorting'' is a problem --- given a sequence of
$n$ numbers, produce a permutation of them in non-decreasing order.

\smallskip
An \term{instance} is one particular input: $\langle 31, 41, 59, 26 \rangle$.

\smallskip
An \term{algorithm} is a finite, unambiguous procedure that produces the
correct output for \emph{every} instance. It is \term{correct} if it halts on
every instance with the right answer --- and ``every'' is doing real work in
that sentence, which is why \chref{invariants} exists.

\smallskip
A \emph{program} is an algorithm written in a language, for a machine. This
book analyses algorithms and measures programs, and is careful about which it
is doing.
\end{definitionbox}

\section{The Model: What Counts as One Step}

To say an algorithm takes ``half a million comparisons'' we need to agree what
a step is. Real machines make that awkward: an addition and a division cost
differently, memory access costs differently depending on the cache, and the
processor may be executing four instructions at once.

So we measure against an idealisation.

\begin{definitionbox}[The Random-Access Machine]
The \term{RAM model} assumes:

\begin{tight}
  \item instructions execute one at a time, with no concurrency;
  \item each \term{primitive operation} --- arithmetic, comparison, assignment,
        reading or writing one memory cell, following one array index --- costs
        one unit of time;
  \item memory is unlimited and every cell costs the same to reach.
\end{tight}

\smallskip
The \term{running time} $T(n)$ of an algorithm is then the number of primitive
operations it performs on an input of size $n$.
\end{definitionbox}

\begin{conceptbox}[Three Lies in That Model, and Why They Are Worth Telling]
Every clause of the definition is false of a real machine, deliberately.

\smallskip
\textbf{Not all operations cost the same.} On a current processor an integer
addition takes about one cycle and an integer division about twenty. The model
prices them equally.

\smallskip
\textbf{Not all memory costs the same.} A read served by the level-1 cache
takes about four cycles; one that goes to main memory takes a few hundred ---
a factor of fifty, which the model ignores entirely. \chref{graphs} is where
this starts to be visible in the measurements.

\smallskip
\textbf{Instructions do not execute one at a time.} Processors reorder,
pipeline and predict branches.

\smallskip
The model is nevertheless the right one, for the reason the next section gives:
every one of those lies is a \emph{constant factor}, and the differences this
course is about are not. A fifty-fold cache penalty is swamped by the
difference between $n$ and $n^2$ before $n$ reaches a thousand. Where the
constant factors do matter --- and they sometimes do --- this book measures
rather than models, which is what the \emph{Try it yourself} boxes are for.
\end{conceptbox}

\section{Why the Question Is Asymptotic}

We do not ask how long an algorithm takes. We ask how the time \emph{grows}
with the input, and we answer with the dominant term, discarding constants.
That sounds like giving up information. It is the opposite: it is the only way
to say something that stays true.

\dsfig{%
\begin{tikzpicture}
\begin{semilogyaxis}[
  width=11.4cm, height=6.4cm,
  xlabel={input size $n$}, ylabel={operations},
  xmin=1, xmax=100, ymin=1, ymax=1e12,
  grid=major, grid style={neutral!20},
  legend pos=outer north east, legend cell align=left,
  legend style={font=\scriptsize, draw=neutral!40},
  tick label style={font=\scriptsize},
  label style={font=\scriptsize},
  every axis plot/.append style={line width=1pt, samples=120}
]
\addplot[greenhl, domain=1:100] {ln(x)/ln(2)};        \addlegendentry{$\lg n$}
\addplot[secondary, domain=1:100] {x};                \addlegendentry{$n$}
\addplot[primary, domain=1:100] {x*ln(x)/ln(2)};      \addlegendentry{$n \lg n$}
\addplot[purplehl, domain=1:100] {x^2};               \addlegendentry{$n^2$}
\addplot[goldhl, domain=1:100] {x^3};                 \addlegendentry{$n^3$}
\addplot[accent, domain=1:40] {2^x};                  \addlegendentry{$2^n$}
\end{semilogyaxis}
\end{tikzpicture}}%
{The six orders of growth this book is mostly about, on a logarithmic vertical
axis --- note that each gridline is a factor of a hundred. At $n = 100$ the
gap between $n\lg n$ and $n^2$ is already a factor of fifteen, and the gap
between $n^2$ and $2^n$ does not fit on the page.}{growth}

\begin{alertbox}[What an exponential actually means]
$2^n$ is not ``slow''. It is a wall.

\smallskip
At $n = 50$, an algorithm doing $2^n$ operations on a machine executing
$10^{10}$ of them a second takes
\[ \frac{2^{50}}{10^{10}} = \frac{1.13 \times 10^{15}}{10^{10}}
   \approx 1.1 \times 10^{5}\ \text{s} \approx 31\ \text{hours}. \]
At $n = 60$ it takes about \textbf{36 days}; at $n = 80$, about a hundred
thousand years. Buying a computer a thousand times faster moves $n = 50$ to
$n \approx 60$ and no further.

\smallskip
That is why \chref{complexity} matters, and why ``we will just use a faster
machine'' is not an answer to it.
\end{alertbox}

\begin{worked}[The Faster Computer Loses]
Computer A executes $10^{10}$ primitive operations a second and runs insertion
sort, which takes about $2n^2$ operations. Computer B executes $10^{7}$ a
second --- a thousand times slower --- and runs merge sort, which takes about
$50 n \lg n$. Sort ten million numbers on each.

\smallskip
\textbf{Step 1 --- Computer A.}
\[ \frac{2 \times (10^{7})^{2}}{10^{10}}
   = \frac{2 \times 10^{14}}{10^{10}}
   = 2 \times 10^{4}\ \text{s} = \mathbf{5\ \text{hours}\ 33\ \text{minutes}} \]

\smallskip
\textbf{Step 2 --- Computer B.} First $\lg(10^{7}) = 7 \lg 10 = 7 \times 3.322
= 23.25$.
\[ \frac{50 \times 10^{7} \times 23.25}{10^{7}}
   = 50 \times 23.25 = 1{,}163\ \text{s}
   = \mathbf{19\ \text{minutes}\ 23\ \text{seconds}} \]

\smallskip
\textbf{Step 3 --- the comparison.} The machine that is \emph{a thousand times
slower} finishes \textbf{17 times sooner}, because its algorithm is better.
The hardware advantage is not small and it is not enough.

\smallskip
\textbf{Step 4 --- where does B overtake A?} Set the two times equal:
\[ \frac{2n^{2}}{10^{10}} = \frac{50 n \lg n}{10^{7}}
   \quad\Longrightarrow\quad 2n = 5 \times 10^{4} \lg n
   \quad\Longrightarrow\quad n = 2.5 \times 10^{4} \lg n. \]
There is no closed form, so iterate from $n = 10^{5}$:
\[ 2.5\!\times\!10^{4}\!\times\!16.6 = 4.15\!\times\!10^{5}
   \;\to\; 4.66\!\times\!10^{5} \;\to\; 4.71\!\times\!10^{5}
   \;\to\; 4.71\!\times\!10^{5}. \]
So A wins below about $n = \mathbf{471{,}000}$ and loses above it --- and loses
by more, the larger the input gets.

\smallskip
\textbf{Step 5 --- what the arithmetic is really saying.} Not that constant
factors are unimportant. A constant of 25 between the two algorithms bought A
nearly half a million elements' worth of advantage, which for many real
problems is the whole range of interest. It is saying that constants buy a
\emph{bounded} advantage and a better order of growth buys an unbounded one,
and that you should know which regime you are in.
\end{worked}

\section{The Claim, Measured}

The worked example is arithmetic about a hypothetical machine. The claim it
rests on --- that a quadratic algorithm's time quadruples when the input
doubles, and an $n\lg n$ one's roughly doubles --- is testable on the machine in
front of you, and ought to be tested.

\begin{conceptbox}[What the ratios should be]
For $T(n) = cn^{2}$,
\[ \frac{T(2n)}{T(n)} = \frac{c(2n)^{2}}{cn^{2}} = 4, \]
independent of $c$ and of $n$. The constant cancels, which is exactly why the
ratio is worth measuring: it tests the \emph{order of growth} without needing
to know anything about the machine.

\smallskip
For $T(n) = cn\lg n$,
\[ \frac{T(2n)}{T(n)} = \frac{2n(\lg n + 1)}{n \lg n} = 2 + \frac{2}{\lg n}, \]
which is a little above 2 and drifts \emph{down} towards 2 as $n$ grows. At
$n = 1000$ it predicts $2 + 2/9.97 = 2.20$.
\end{conceptbox}

\begin{worked}[The Prediction Against the Measurement]
Running \texttt{code/ch01\_growth.cpp} on the machine described in Appendix~A
(MSVC 19.44, \texttt{/O2}), median of five runs, random permutations:

\rowcolors{2}{white}{lightbg}
\begin{center}
\begin{tabular}{@{}rrrrr@{}}
\toprule
\rowcolor{primary!15}
$n$ & insertion (ms) & ratio & merge (ms) & ratio \\
\midrule
1{,}000  &   0.15 & ---  & 0.06 & ---  \\
2{,}000  &   0.58 & 3.87 & 0.14 & 2.22 \\
4{,}000  &   2.41 & 4.13 & 0.30 & 2.20 \\
8{,}000  &   9.35 & 3.88 & 0.65 & 2.14 \\
16{,}000 &  37.08 & 3.97 & 1.51 & 2.32 \\
32{,}000 & 145.86 & 3.93 & 2.94 & 1.96 \\
\bottomrule
\end{tabular}
\end{center}

\smallskip
\textbf{Insertion sort.} Predicted 4; measured 3.87, 4.13, 3.88, 3.97, 3.93.
Mean 3.96.

\smallskip
\textbf{Merge sort.} Predicted $2 + 2/\lg n$ --- 2.20 at $n = 1000$, falling to
2.14 by $n = 8000$. Measured 2.22, 2.20, 2.14, 2.32, 1.96. The first three
agree to within a per cent; the last two straddle the prediction.

\smallskip
\textbf{Read the last column honestly.} The 2.32 and the 1.96 are not the
theory failing. They are two measurements of a quantity near 2.15 with a few
per cent of noise, and a ratio magnifies noise in both terms. The way to tell
noise from a real effect is to look at the \emph{pattern}: insertion's five
ratios scatter about 4 and merge's about 2.2, and neither drifts.

\smallskip
\textbf{And notice what the measurement cost.} At $n = 32{,}000$ insertion
sort took 146 ms and merge sort 2.9 ms --- a factor of 50. At $n = 1{,}000$ the
factor was 2.5. The gap is not fixed; it widens with every doubling, and that
is the whole content of ``asymptotically faster''.
\end{worked}

\begin{alertbox}[The asymptotics are not the whole story, and this book says so twice]
The same program, run on small inputs, gives this:

\rowcolors{2}{white}{lightbg}
\begin{center}
\begin{tabular}{@{}rrrl@{}}
\toprule
\rowcolor{primary!15}
$n$ & insertion ($\mu$s) & merge ($\mu$s) & \\
\midrule
 20 & 0.20 & 0.30 & insertion wins \\
 40 & 0.50 & 0.60 & insertion wins \\
 60 & 0.70 & 1.00 & insertion wins \\
 70 & 1.20 & 1.10 & merge wins \\
100 & 2.20 & 1.60 & merge wins \\
200 & 7.40 & 4.60 & merge wins \\
\bottomrule
\end{tabular}
\end{center}

\smallskip
The crossover is at about $n = 65$ on this machine. Below it the
``worse'' algorithm is faster, every time, because merge sort's recursion and
its auxiliary array cost more than insertion sort's whole quadratic loop on
sixty elements.

\smallskip
This is not a curiosity. It is why every production sort in the world --- the
C++ standard library's included --- switches to insertion sort below a
threshold of roughly this size. \chref{quicksort} returns to it.
\end{alertbox}

\section{Space, and the Trade Against Time}

Time is not the only resource. \term{Space complexity} counts the memory an
algorithm needs beyond its input, and the two are often tradeable: a lookup
table turns a computation into a memory access, and \chref{dpone}'s whole
technique is spending memory to avoid recomputation.

\begin{conceptbox}[Which Resource Binds?]
For most of the twentieth century, memory was the scarce one. It is now
usually not: a machine with 32 GB can hold an array of four billion integers,
and most problems that fit in a lecture do not come close.

\smallskip
What has replaced it is \emph{memory locality}. An algorithm using twice the
memory but touching it in order can comfortably beat one using half as much in
a scattered pattern, because of the cache the RAM model promised to ignore.
\chref{graphs} measures exactly this.

\smallskip
So: analyse space because the exam asks and because occasionally it binds; but
when two algorithms have the same order of growth and different speeds, look at
the access pattern before the byte count.
\end{conceptbox}

\begin{pitfall}
\begin{tight}
  \item \textbf{Confusing the algorithm with the program.} ``My sort is slow''
        is usually a statement about an implementation; ``insertion sort is
        quadratic'' is a statement about an idea. Fixing the first will not
        change the second.
  \item \textbf{Timing an unoptimised build.} A debug build can be five to ten
        times slower and --- worse --- slower by \emph{different} factors for
        different code, so the comparison is meaningless. Appendix~A insists
        on \texttt{/O2}.
  \item \textbf{Timing once.} A single run contains whatever else the machine
        was doing. Every measurement in this book is the median of several.
  \item \textbf{Believing the asymptotics at small $n$.} The crossover above is
        at $n = 65$. If your inputs are always smaller than that, the better
        order of growth is buying you nothing.
  \item \textbf{Believing the constants matter at large $n$.} The symmetric
        error. At $n = 32{,}000$ the ``unimportant'' factor of 50 was entirely
        due to the order of growth.
  \item \textbf{Reporting a mean of timings.} One scheduler interruption drags
        a mean a long way. Use the median, and say how many runs.
\end{tight}
\end{pitfall}

\begin{lab}[Growth, Measured]
Reproduces both tables above on your own machine. Appendix~A sets up the
compiler. Expect twenty minutes, most of it reading the output.

\begin{lstlisting}[language=C++]
// ch01_growth.cpp -- what "the question is asymptotic" looks like.
// Build: cl /O2 /EHsc /std:c++17 ch01_growth.cpp      (see Appendix A)

// --- 1. The two algorithms, written plainly -------------------------
static void insertion_sort(std::vector<int>& a) {
    for (size_t j = 1; j < a.size(); ++j) {
        int key = a[j];
        size_t i = j;
        while (i > 0 && a[i - 1] > key) { a[i] = a[i - 1]; --i; }
        a[i] = key;
    }
}

static void merge_sort(std::vector<int>& a, std::vector<int>& buf,
                       size_t lo, size_t hi) {
    if (hi - lo < 2) return;
    size_t mid = lo + (hi - lo) / 2;
    merge_sort(a, buf, lo, mid);
    merge_sort(a, buf, mid, hi);
    size_t i = lo, j = mid, k = lo;
    while (i < mid && j < hi) buf[k++] = (a[i] <= a[j]) ? a[i++] : a[j++];
    while (i < mid) buf[k++] = a[i++];
    while (j < hi)  buf[k++] = a[j++];
    std::copy(buf.begin() + lo, buf.begin() + hi, a.begin() + lo);
}

// --- 2. Timing that is honest ---------------------------------------
// The MEDIAN of several runs: a mean is dragged by one scheduler
// hiccup, and the minimum flatters the cache.
template <class F>
static double time_ms(F&& f, int reps) {
    std::vector<double> t;
    for (int r = 0; r < reps; ++r) {
        auto t0 = std::chrono::steady_clock::now();
        f();
        auto t1 = std::chrono::steady_clock::now();
        t.push_back(
            std::chrono::duration<double, std::milli>(t1 - t0).count());
    }
    std::sort(t.begin(), t.end());
    return t[t.size() / 2];
}

// --- 3. Double n each time, and watch the ratios --------------------
// Predicted: 4 for insertion sort, 2 + 2/lg(n) for merge sort.
for (size_t n = 1000; n <= 32000; n *= 2) {
    auto base = random_input(n, 12345);
    std::vector<int> buf(n);
    double ti = time_ms([&]{ auto a = base; insertion_sort(a); }, 5);
    double tm = time_ms([&]{ auto a = base;
                             merge_sort(a, buf, 0, n); }, 5);
    printf("%8zu %12.2f %7.2f %12.2f %7.2f\n",
           n, ti, ti / prev_i, tm, tm / prev_m);
    prev_i = ti; prev_m = tm;
}

// --- 4. Now find where they cross -----------------------------------
// The other half of the point: below the crossover the "worse"
// algorithm is faster, and every production sort knows it.
for (size_t n = 10; n <= 200; n += 10) {
    /* same two timings, many more repetitions because the times
       are microseconds; print both and mark which one won */
}
\end{lstlisting}

\textbf{What to notice.} Your absolute times will differ from the book's --- a
different processor, a different compiler, a different afternoon. The
\emph{ratios} should not. If your insertion-sort column is not near 4, something
is wrong with the measurement, not with the theory: check that you built with
optimisation on, and that nothing else was running.
\end{lab}

\begin{checkpoint}
\begin{tightnum}
  \item An algorithm takes 3 seconds on an input of size 1{,}000 and 48 seconds
        on an input of size 4{,}000. What is its order of growth, and how did
        you decide?
  \item Computer A runs $n^{3}$ operations at $10^{9}$ per second; computer B
        runs $100 n^{2}$ at $10^{7}$ per second. At $n = 10^{4}$, which
        finishes first, and by what factor?
  \item The RAM model prices a cache hit and a main-memory access the same,
        although they differ by a factor of about fifty. Give one situation in
        which that lie is harmless and one in which it is not.
\end{tightnum}
\end{checkpoint}

\begin{chaptersummary}
\begin{tight}
  \item A problem is a specification, an instance is one input, an algorithm is
        a procedure correct on \emph{every} instance. A program implements an
        algorithm; this book analyses the second and measures the first.
  \item The RAM model costs every primitive operation at one unit. All three of
        its assumptions are false of real hardware, and all three are constant
        factors --- which is why the model is still the right one.
  \item We ask how running time \emph{grows}, not what it is, because the
        growth is the part that survives a change of machine.
  \item Constants buy a bounded advantage; a better order of growth buys an
        unbounded one. A computer 1{,}000 times faster still lost by a factor
        of 17, and the crossover was at $n \approx 471{,}000$.
  \item Doubling the input multiplies a quadratic time by 4 and an $n \lg n$
        time by $2 + 2/\lg n$. Measured on this machine: 3.96 and 2.17.
  \item Below $n \approx 65$ insertion sort beat merge sort on the same
        machine. Asymptotics describe the limit, and the limit is not
        everywhere.
  \item Measure with optimisation on, several times, and report the median.
\end{tight}
\end{chaptersummary}

\begin{reviewq}
  \item \textbf{(MCQ)} The RAM model assumes: (a)~memory access is free;
        (b)~every primitive operation costs one unit of time; (c)~instructions
        execute in parallel; (d)~arithmetic is exact to arbitrary precision.
  \item \textbf{(MCQ)} If $T(n) = cn^{2}$, then $T(2n)/T(n)$ equals:
        (a)~2; (b)~$2 + 2/\lg n$; (c)~4; (d)~it depends on $c$.
  \item \textbf{(Short)} Explain why discarding constant factors makes an
        analysis \emph{more} useful rather than less.
  \item \textbf{(Short)} Give the three assumptions of the RAM model and, for
        each, one way real hardware violates it.
  \item \textbf{(Short)} A colleague says ``we do not need a better algorithm,
        we will buy a faster server''. Give the shortest argument that settles
        the question, and say what information you would need first.
  \item \textbf{(Applied)} Algorithm P takes $0.5 n^{2}$ operations and
        algorithm Q takes $200 n \lg n$. Find the input size at which Q
        overtakes P, showing the iteration. Then say what you would recommend
        to a team whose inputs never exceed 2{,}000.
  \item \textbf{(Applied)} Design a measurement that would distinguish an
        $\bigO{n\lg n}$ algorithm from an $\bigO{n}$ one, given only a compiled
        binary you can time. State the sizes you would use, how many
        repetitions, what you would compute, and what result would make you
        say the evidence is inconclusive.
\end{reviewq}
